% Appendix: designs and full tables of the mechanism experiments. Numbers recomputed from the prediction files in
% notes/paper/mechanism_check.md; materials, code and cards in the results repository (notes/2026-11-17-mechanism).

\section{Mechanism Experiments}
\label{app:mechanism}

\paragraph{Arithmetic chains.} Each question lists five chains of variable assignments, interleaved in random order.
The chains share a root value and a sequence of operation types and differ only in their constants: the first hop
adds or subtracts a constant, later hops add, subtract, double or triple, with constants from 2 to 9 and values kept
between 1 and 99. The question asks for the variable $k$ hops from the root, and the five options are the five chains'
values at that depth; an item is redrawn when two options coincide, an option appears in the text, or one option is
twice or three times another. The test set has 150 items for each $k$ from 1 to 6 and 100 for $k = 7, 8$; chance is
20\%. Training examples come from the same generator with another seed and cover $k = 1$ to 6, so seven and eight hops
are extrapolation. The answer-only weights of Figure~\ref{fig:mechanism}a are one-epoch GSM8K runs (seeds 16--20) evaluated
on the chain sets without any chain training; the chain-trained runs add 27,000 chain examples (26,941 distinct value
sets) to the 5,923 GSM8K problems in one epoch (1,029 updates instead of 186). At four hops and more, a reader that
gets every constant but one right still scores high on this option design, and few items require the full chain; we
therefore regard three hops as the cleanest test of composition. Smaller data budgets (3,000 and 9,000 chain examples,
with and without repetition) give intermediate curves, and a two-stage variant that keeps the chain capacity pays a
similar GSM8K cost as a one-hop control trained the same way, so the cost comes with the extra training rather than
with composition.

\paragraph{The GSM8K transfer test of chain training.} For the GSM8K test problems whose critical path has at least
four serial steps, the test shows the true steps up to the point where two serial steps remain (the registered cell
``rc2''), on the 249 problems where no shown value equals an option. Compared with one-epoch answer-only training at
the same five seeds, the chain-trained models are $7.6$ points less accurate (per seed $-4.4$, $-2.4$, $-7.6$, $-12.9$,
$-10.4$; seed-level interval $[-12.8, -2.3]$, two-level $[-13.3, -2.0]$). Against a second set of answer-only runs of
the same configuration, the difference is $-6.2$ (two-level $[-11.6, -1.0]$), and against that second set on the
plain GSM8K test with the original options it is $-1.6$ (two-level $[-3.5, +0.2]$).

\paragraph{Hint test: items and conditions.} The 1,129 items are the GSM8K test problems whose worked solution has at
least two steps, where a step is one binary operation taken from the calculator notes of the Socratic solution; the
options are the cue-free ones (the four wrong options are answers of other problems). A ``Worked steps:'' block
between the question and the options shows the first $n - r$ true steps, so that $r$ steps are left to compute; the
sets shrink with $r$ (1,129, 877, 576, 332, 158 and 75 items for $r = 1$ to 6). The controls use the same number of
shown lines as the true condition at the same $r$: the problem's own sentences (length control), the first steps of
another problem's solution, a consistently wrong chain in which every result is off by at least 1.5 and later
operands follow the wrong values, and the true operands with results changed by a small amount without propagation.
A ``full'' set shows every step, including the answer, as a ceiling check.

\paragraph{Hint test: models and registration.} A pilot on saved four-epoch weights (seeds 12--14) was used only to
decide whether to run the test. The registered test trained fresh runs at seeds 16--18 (answer-only for one and two
epochs, \tf{} and the permuted-nodes control for two epochs, data and settings of Section~\ref{sec:gsm8k}) and evaluated the
hint sets at the end of training. Its primary statistic is the gain from the last-step-left condition over no steps
on the 440 problems with at least four steps where no shown value equals an option, for the answer-only model selected
on validation (one epoch); the criterion (mean at least 15 points, every seed at least 10, seed-level $t$ above the
critical value) is met with $+36.4$ (per seed $+38.2$, $+34.1$, $+37.0$). The \tf{} runs narrowly fail the registered
length control, which bounds the drop caused by the problem's own text at 3 points: with two lines of their own text
shown, they are 4.1 points less accurate than with none (both intervals below zero), against at most 1.4 points for
the other methods.

\begin{table}[h]
\centering
\scriptsize
\setlength{\tabcolsep}{3pt}
\caption{Hint test accuracy (\%) by condition, seed mean over seeds 16--18, each set on its own items. True $r$: the
first $n-r$ steps shown. Own text, other problem, wrong chain and near-miss results: controls with the same number of
shown lines as true $r$.}
\label{tab:hintfull}
\resizebox{\linewidth}{!}{%
\begin{tabular}{lcccccccccccccc}
\toprule
 & None & All & \multicolumn{4}{c}{True $r$} & \multicolumn{2}{c}{Own text $r$} & \multicolumn{2}{c}{Other problem $r$} & \multicolumn{2}{c}{Wrong chain $r$} & \multicolumn{2}{c}{Near-miss $r$} \\
\cmidrule(lr){4-7}\cmidrule(lr){8-9}\cmidrule(lr){10-11}\cmidrule(lr){12-13}\cmidrule(lr){14-15}
Method & & & 1 & 2 & 3 & 4 & 1 & 2 & 1 & 2 & 1 & 2 & 1 & 2 \\
\midrule
Answer-only, 1 epoch & 60.6 & 95.6 & 85.9 & 67.0 & 53.0 & 45.8 & 60.8 & 52.3 & 53.2 & 46.4 & 51.3 & 48.0 & 62.0 & 53.0 \\
Answer-only, 2 epochs & 59.0 & 97.6 & 83.4 & 66.0 & 51.2 & 45.8 & 58.8 & 50.5 & 52.0 & 46.5 & 53.3 & 47.7 & 60.4 & 50.8 \\
\tf{}, 2 epochs & 60.8 & 96.1 & 87.9 & 65.8 & 51.2 & 43.8 & 58.4 & 49.4 & 50.3 & 43.6 & 50.5 & 46.3 & 61.0 & 51.9 \\
Permuted nodes, 2 epochs & 58.6 & 96.9 & 83.9 & 64.7 & 50.5 & 42.8 & 57.9 & 50.2 & 48.9 & 43.5 & 51.7 & 47.0 & 59.8 & 51.0 \\
\bottomrule
\end{tabular}}
\end{table}

\begin{table}[h]
\centering
\small
\caption{Hint test on the 332 problems with at least five steps (fixed items): accuracy (\%) by the number of steps
left, and $D$ = one step left minus four steps left (points, two-level 95\% interval).}
\label{tab:hintcurve}
\begin{tabular}{lccccccc}
\toprule
Method & None & 4 left & 3 left & 2 left & 1 left & All shown & $D$ \\
\midrule
Answer-only, 1 epoch & 45.7 & 45.8 & 54.2 & 62.9 & 81.7 & 95.7 & $+35.9$ [$+30.3$, $+41.6$] \\
Answer-only, 2 epochs & 44.7 & 45.8 & 49.1 & 61.2 & 76.9 & 97.3 & $+31.1$ [$+25.3$, $+36.7$] \\
\tf{}, 2 epochs & 43.8 & 43.8 & 50.4 & 60.4 & 82.0 & 96.1 & $+38.3$ [$+32.1$, $+44.3$] \\
Permuted nodes, 2 epochs & 41.7 & 42.8 & 49.1 & 58.3 & 77.7 & 97.8 & $+34.9$ [$+28.5$, $+41.3$] \\
\bottomrule
\end{tabular}
\end{table}

\paragraph{Tree Forcing against answer-only training in the hint test.} Paired by seed, the gain of the \tf{} runs
from the last-step-left condition over no steps is $+1.8$ points larger than that of one-epoch answer-only training
(per seed $+3.2$, $+1.7$, $+0.6$; seed-level interval $[-1.4, +5.0]$, two-level $[-1.0, +4.7]$) and $+1.9$ larger than
that of the permuted-nodes control (two-level $[-1.5, +5.5]$); neither is significant. In the last-step-left condition
itself, the \tf{} runs are 4.5 and 4.1 points more accurate than the two-epoch answer-only and permuted-nodes runs
(both significant), but not than one-epoch answer-only training ($+2.0$, seed-level interval $[-0.5, +4.5]$).

\begin{table}[h]
\centering
\small
\caption{One-pass accuracy (\%) on the arithmetic chains by hop count (seed mean; five seeds except the untrained
model). Chance is 20\%; seven and eight hops were not trained.}
\label{tab:khop}
\begin{tabular}{lcccccccc}
\toprule
Model & 1 & 2 & 3 & 4 & 5 & 6 & 7 & 8 \\
\midrule
Untrained & 29.3 & 18.7 & 18.0 & 20.7 & 14.7 & 16.0 & 18.0 & 24.0 \\
Answer-only (GSM8K, 1 epoch) & 93.6 & 42.0 & 27.9 & 23.1 & 22.7 & 21.1 & 23.8 & 25.0 \\
+ 27,000 chain examples & 100.0 & 99.9 & 90.7 & 87.3 & 76.3 & 67.2 & 55.0 & 47.4 \\
\bottomrule
\end{tabular}
\end{table}
